\partstart{Phân tích mã nguồn và chức năng}
\subsection{Tổ chức repository và luồng phụ thuộc}
Repository tách mã thực thi trong \code{src}, cấu hình trong \code{configs},
notebook diễn giải trong \code{notebooks}, kiểm thử trong \code{tests},
kết quả tổng hợp trong \code{results} và báo cáo trong \code{report}.
Tách biệt này giúp người đọc phân biệt dữ liệu đầu vào, chương trình tạo
kết quả và tài liệu phân tích. Không dùng ảnh chụp notebook làm nguồn duy
nhất của metric khi CSV/JSON đã lưu có thể đối chiếu trực tiếp.

\begin{table}[!htbp]\centering\small
\caption{Cấu trúc và trách nhiệm của các thành phần chính.}\label{tab:repo}
\begin{tabular}{p{3.4cm}p{9.9cm}}\toprule
Đường dẫn & Vai trò và sản phẩm\\\midrule
\code{data/raw} & CSV nguồn cục bộ, có hướng dẫn tải; dữ liệu lớn không đưa vào Git.\\
\raggedright\code{configs/}\newline\code{default.json} & Đường dẫn nguồn, seed, ngân sách dòng, test size, số fold và số tiến trình.\\
\code{src} & Đọc dữ liệu, pipeline, runner, kiểm toán, đánh giá, xuất báo cáo.\\
\code{notebooks} & Bốn notebook diễn giải EDA, tiền xử lý, thí nghiệm và phân tích lỗi.\\
\code{results/runs} & Từng lần chạy: manifest, split, mô hình, bảng và hình; không phải toàn bộ đều công bố.\\
\code{results/tables} & Bảng tổng hợp đã lọc, metric, audit, sampling và phân tích lỗi.\\
\code{results/figures} & Hình EDA và đánh giá dùng trong báo cáo.\\
\code{tests} và \code{docs} & Kiểm thử phần mềm; yêu cầu học phần, bằng chứng và hướng dẫn.\\\bottomrule
\end{tabular}

\end{table}
Bảng~\ref{tab:repo} cũng cho thấy ranh giới công bố: split chứa chỉ số
nguồn và mô hình nhị phân được giữ cục bộ, còn bảng aggregate có thể được
đưa vào báo cáo. Các tệp lịch sử nằm ở thư mục riêng; không trộn metric từ
hai protocol mà không chỉ rõ nguồn. Run công bố hiện tại được xác định
bằng manifest và hash dữ liệu.

\subsection{Module data.py: đọc, kiểm tra và chia nhóm}
Hàm \code{read_csv} lần lượt thử UTF-8 có BOM và CP1252; chỉ chuyển encoding
khi gặp lỗi giải mã. Nếu không đọc được, chương trình báo lỗi thay vì dùng
cơ chế bỏ qua byte. Hàm \code{sha256} đọc từng khối byte để tạo dấu vân tay
tệp, giúp kiểm tra cùng tên CSV có thực sự cùng nội dung hay không.

Các hằng \code{NUM}, \code{CAT}, \code{TARGET}, \code{GROUP} và \code{DATE}
tạo hợp đồng chung giữa module dữ liệu và pipeline. Nhờ vậy, danh sách
cột sử dụng không bị viết lại khác nhau giữa notebook và runner. Đầu vào
thiếu cột bắt buộc sẽ dừng ngay trong \code{validate}; ngày lỗi cũng được
phát hiện trước khi trích đặc trưng lịch.

Đoạn mã sau có trong \code{src/data.py}, thể hiện kiểm tra điều kiện
không giao đơn giữa hai tập sau khi bộ chia trả chỉ số:
\begin{verbatim}
if set(df.iloc[tr][GROUP]) & set(df.iloc[te][GROUP]):
    raise AssertionError('Order leakage')
\end{verbatim}
Kiểm tra trực tiếp làm rõ điều kiện cần bảo vệ, đồng thời giúp phát hiện
sai sót khi thay bộ chia. Hàm \code{protocol_split} gói lấy mẫu nguyên đơn
và split để notebook có thể dùng cùng thứ tự thao tác với runner. Tái lập
còn cần thứ tự dòng nguồn và phiên bản thư viện, không chỉ cùng seed.

\code{feature_audit} duyệt mọi cột, ghi keep/drop, lý do và độ chắc chắn.
Điều này giúp truy vết vì sao không dùng một trường có vẻ dự đoán tốt.
Chẳng hạn, loại trạng thái giao là để giữ đúng mục tiêu dự báo sớm, trong
khi loại trường không nằm trong allowlist có thể là lựa chọn bảo thủ do
ngữ nghĩa chưa đủ rõ. Hai lý do cần được phân biệt trong báo cáo.

\subsection{Module pipeline.py: đặc trưng và phép biến đổi}
\code{SafeFeatures} là transformer theo giao diện estimator của
scikit-learn. Phương thức fit kiểm tra tên luồng A/B; transform tạo bản sao
chỉ chứa cột cho phép. Luồng B dùng ngày tạo đơn để bổ sung lịch. Không có
bước sử dụng nhãn để tạo đặc trưng, và các cột kết quả không đi qua
transformer này.

\begin{verbatim}
if self.flow == 'B':
    dt = pd.to_datetime(X[DATE], errors='raise')
    out['order_month'] = dt.dt.month
    out['order_weekday'] = dt.dt.dayofweek
    out['order_hour'] = dt.dt.hour
\end{verbatim}
Đoạn mã cho thấy B không sử dụng thời gian giao thực tế. Tuy nhiên, trích
đặc trưng lịch đúng cú pháp không tự xác nhận múi giờ hay ý nghĩa vận hành
của timestamp. Khi dùng dữ liệu ngoài, cần kiểm tra hai khía cạnh đó trước
khi áp dụng cùng quy tắc. Báo cáo không suy diễn rằng giờ tạo đơn là giờ
địa phương ở mọi quốc gia.

Hàm \code{build} ghép ba bước có tên features, preprocess và model.
\code{ColumnTransformer} chuyển nhánh số qua median/scaling, nhánh phân
loại qua mode/one-hot. Tên tham số tìm kiếm có tiền tố \code{model__},
giúp thay cấu hình mô hình mà vẫn giữ phép tiền xử lý trong mỗi lần fit.
Đầu ra thưa được dùng cho cả hai thuật toán để thống nhất cấu trúc.

Hàm \code{grid} trả tám cấu hình K-NN hoặc 12 cấu hình cây. Baseline và
Dummy bỏ qua grid, nhưng vẫn được cross-validation trên cùng fold. Đây
là điểm cần lưu ý khi đọc thời gian: thời gian search/refit của mô hình
tuned gồm nhiều lần fit, không thể so trực tiếp như thời gian một lần
huấn luyện cây baseline.

\subsection{Module run.py: điều phối thí nghiệm}
Runner đọc cấu hình, xác thực nguồn và tạo thư mục run có timestamp UTC.
Manifest ban đầu mang trạng thái RUNNING; nếu có ngoại lệ, trạng thái đổi
thành FAILED và lưu lỗi. Chỉ sau khi hoàn tất bảng, hình và mô hình, runner
ghi COMPLETED. Cách ghi này giúp tránh xuất bản một lần chạy dở như kết
quả cuối cùng.

Trình tự quan trọng là: tạo fold trên train; fit sáu ứng viên; chọn winner
theo CV; lưu manifest; chạy chẩn đoán ablation trên train; sau đó mới tính
metric test. Đoạn mã thực hiện khóa người thắng trước test là:
\begin{verbatim}
winner=max(candidates,key=lambda x:x[1]['cv_f1'])[1]['name']
manifest['selected_by_cv']=winner; persist()
\end{verbatim}
Do đó, tree B được gọi là mô hình chọn theo CV, không phải mô hình được
chọn sau khi xem F1 test. Phân tích lại lỗi test trong báo cáo không tạo
thêm một lần xác nhận độc lập và không được dùng để chọn lại tham số.

Mỗi mô hình được lưu bằng joblib, nạp lại và kiểm tra dự đoán giống bản
trước khi lưu. Runner ghi classification report, ma trận nhầm lẫn, thời
gian, metric xác suất và F1 train. Với mô hình chọn, ví dụ lỗi chỉ giữ các
trường nghiệp vụ cho phép, không chứa mã đơn, tên, địa chỉ hay timestamp
chính xác. Việc lọc này giảm thông tin cá nhân trong artifact công bố.

\subsection{Module evaluation.py: ba chẩn đoán bổ sung}
\code{representativeness} so sánh phân phối danh mục, mean/median trường
số và kích thước đơn giữa toàn nguồn với mẫu. Các chênh lệch là thống kê
mô tả. Hàm không thực hiện kiểm định đầy đủ cho mọi quan hệ hay chứng minh
mẫu không chệch. Có cùng phân phối biên vẫn có thể khác quan hệ điều kiện
giữa đặc trưng và nhãn.

\code{order_metrics} nhóm xác suất theo mã đơn, kiểm tra nhãn nhất quán,
lấy trung bình rồi dùng ngưỡng 0,5. Nếu có nhãn mâu thuẫn, hàm ghi BLOCKED
thay vì tùy ý chọn nhãn đầu tiên. Bước này làm rõ tiền đề để mỗi đơn có
trọng số bằng nhau, khác với metric mặc định tính theo số dòng.

\code{time_ablation} lấy tham số của ứng viên A đã chọn, tạo lại A và B
với cùng tham số, rồi dùng cùng fold train để đánh giá. Tên cột
\code{test_f1} là quy ước trả về của \code{cross_validate}; trong tệp
ablation nó chỉ điểm validation của CV, không phải test holdout.
Phân biệt này ngăn so sánh nhầm 0,6904 của ablation cây B với 0,7098 của
test cây B có cấu hình được chọn riêng.

\subsection{Module audit.py và reporting.py}
\code{audit.py} xuất từ điển dữ liệu với kiểu, số thiếu, cardinality và mô
tả từ tệp nguồn. Nó còn xuất thống kê số, số outlier IQR, số đơn nhiều dòng
và độ nhất quán nhãn. Hàm plots nhận train sau split, tạo tám hình khảo
sát. Biểu đồ vùng biểu diễn tỷ lệ giao trễ, kèm bảng số lượng vùng; không
được nhầm chiều cao cột với số lượng đơn ở vùng đó.

\code{reporting.py} chỉ nhận run COMPLETED, đọc comparison và manifest để
chọn đúng hàng của winner. Hàm tạo summary tiếng Việt, bảng LaTeX và bản
nguồn tổng hợp tự động. Nó sao chép bảng công bố nhưng loại split.csv;
không sao chép model hoặc dữ liệu thô. Bản báo cáo hoàn chỉnh này dùng
nguồn biên tập ở \code{report/sections}, còn các nguồn tự động cũ được
giữ để truy vết, không chèn lặp toàn bộ hình và kết quả.

\subsection{Notebook, môi trường và thực thi}
Bốn notebook lần lượt là \code{01_eda.ipynb},
\code{02_preprocessing.ipynb}, \code{03_model_experiments.ipynb} và
\code{04_error_analysis.ipynb}. Chúng sử dụng module chung để hạn chế
sai khác giữa quá trình khám phá và runner. Notebook giúp trình bày từng
bước nhưng một ô hiển thị thành công không thay thế manifest, hash và
bảng metric của lần chạy được sử dụng.

Cấu hình mặc định đặt seed 42, max rows 6.000, test size 0,2, CV ba fold
và n jobs 1. Người tái lập cần tải đúng hai CSV vào đường dẫn trong config,
cài requirements rồi chạy \code{py -m src.run --config configs/default.json}
từ gốc repository. Đây là lệnh tái lập được mô tả, không phải lệnh đã được
chạy lại để soạn báo cáo này. Kết quả lịch sử vẫn là cơ sở số liệu duy nhất.

Manifest ghi Python 3.13.6, NumPy 2.4.6, pandas 3.0.5, scikit-learn 1.9.0,
matplotlib 3.10.9, seaborn 0.13.2 và joblib 1.6.0. Đây là phiên bản của lần
chạy đã lưu, không suy ra phiên bản của mọi môi trường đọc báo cáo. Cùng
seed nhưng khác phiên bản hoặc thứ tự dòng vẫn có thể khác kết quả.

\subsection{Kiểm thử đã có và giới hạn bằng chứng}
\begin{table}[!htbp]\centering\small
\caption{Phạm vi kiểm thử phần mềm trong repository.}\label{tab:tests}
\begin{tabular}{p{4.2cm}p{9.1cm}}\toprule
Nhóm kiểm tra & Điều kiện được kiểm soát\\\midrule
Schema và dữ liệu lỗi & Thiếu nhãn, nhóm, ngày; kiểu cột số; hợp đồng dữ liệu thật nếu có.\\
Nhóm và protocol & Không giao đơn train/test, không giao nhóm CV, cùng protocol cho kết quả lặp.\\
Pipeline và rò rỉ & Danh mục lạ, missing, hình dạng đầu ra, train-only imputation; thay cột cấm không đổi dự đoán.\\
Lưu trữ và đánh giá & Dự đoán sau nạp model; trọng số đơn bằng nhau; chặn nhãn mâu thuẫn.\\
Smoke tổng hợp & Runner đi hết quy trình trên fixture nhỏ; không phải thực nghiệm DataCo.\\\bottomrule
\end{tabular}

\end{table}
Hồ sơ \code{REPAIR_VERIFICATION.md} ghi lần kiểm thử workspace 14 passed,
6 warnings và clean export 13 passed, một skip dự kiến, 6 warnings. Tệp
\code{pytest_output.txt} cũ ghi 12 passed; các con số khác nhau thuộc các
mốc kiểm tra khác nhau, không được gọi tất cả là cùng một lần chạy.
Bảng~\ref{tab:tests} mô tả nội dung test hiện hữu; báo cáo không coi fixture
tổng hợp là bằng chứng hiệu quả mô hình và không tuyên bố chạy mới bộ test
ML trong lần biên tập này.

Kiểm thử phần mềm không xác nhận ngữ nghĩa nghiệp vụ, tính ổn định của
dữ liệu tương lai hoặc lợi ích can thiệp. Những vấn đề này cần đánh giá
độc lập ngoài unit test.
